\chapter*{\zihao{-3}\heiti{摘~~~~要}}
基于物理的图形仿真技术在游戏和电影特效中被大量使用，流体同弹塑性固体的高精度交互仿真已经成为研究的热点。物质点法（MPM）是一种可以实现流体、弹性体、刚体等多种材料相互耦合的仿真方法，在诸多应用领域已取得显著成果，但它在算法层面仍存在固有问题，例如物质点法在模拟流固耦合时会产生数值耗散，不同材质的物体在耦合边界处会出现网格边长大小的空隙。另一方面，影视和游戏行业对物理仿真解算时间效率的需求日益提高，因此GPU优化算法在物理仿真中愈发重要。针对上述问题，本文基于物质点法，提出一套融合自适应网格、支持非连续速度场的流固耦合仿真方法，并进一步给出GPU优化方案。

在算法方面，针对传统物质点法在统一背景网格上存在数值粘性的局限，本文通过修改物质点法解算流程，引入分解兼容性仿射粒子网格（DC-APIC）积分器，在切向方向强制实施自由滑移边界条件，在法向方向实现分离条件，并设计基于相场梯度的分割方法，经多种复杂的流体与弹塑性固体相互作用仿真，验证了该方法的有效性。针对DC-APIC积分器在固体尖锐边缘处数值不稳定的问题，以及传统物质点法处理物体碰撞边界存在非物理空隙的现象，本文提出融合DC-APIC积分器的自适应网格算法，为适应流固耦合条件下动态划分网格的需求，本文使用基于相场梯度的粒子层级来驱动网格划分，实验表明，该方法在不增加过多额外开销的同时，有效保留不同尺度的碰撞细节，能够实现高精度的流固界面仿真。在并行优化方面，本文基于GPU-MPM算法框架，提出了面向非连续速度场的GPU加速方案，设计了一种针对非连续数据的线程束写入优化算法，在高精度大规模场景仿真测试中，本文的GPU优化算法的计算效率相较于SIMD加速的CPU端物质点法达到了17.7到25.7倍。

综合DC-APIC粒子网格传输方法、自适应物质点法和GPU优化策略，本文的流固耦合仿真框架可以高效地仿真常见流体与弹塑性固体的交互场景，并在流固交界面保留高分辨率效果和流固两相自然分离的物理正确性。同时，本文的框架仅依赖单个背景网格，不需要复杂的几何边界约束，实现简单，适合用于大规模流固场景仿真的特效制作。


{\sihao{\textbf{关键词：}}} 计算机图形学，物理仿真，物质点法，流固耦合，自适应性网格
\addcontentsline{toc}{chapter}{摘要}